% Cheese demand, retail conduct and a ban on store-cut cheese
% LaTeX version of basics.md. Compiles with pdflatex (or Overleaf).
\documentclass[11pt]{article}

\usepackage[margin=1in]{geometry}
\usepackage[T1]{fontenc}
\usepackage[utf8]{inputenc}
\usepackage{lmodern}
\usepackage{amsmath,amssymb}
\usepackage{booktabs}
\usepackage{tabularx}
\usepackage{array}
\usepackage{enumitem}
\usepackage[authoryear,round]{natbib}
\usepackage[hidelinks]{hyperref}

\setlist{itemsep=2pt, topsep=4pt}
\newcolumntype{L}{>{\raggedright\arraybackslash}X}
\newcolumntype{C}{>{\centering\arraybackslash}X}
\renewcommand{\arraystretch}{1.15}

\newcommand{\code}[1]{\texttt{#1}}
\newcommand{\model}[1]{\textsf{#1}}

\title{Cheese Demand, Retail Conduct and a Ban on Store-Cut Cheese:\\
Dominick's, 1989--1997}
\author{}
\date{}

\begin{document}
\maketitle

This note is the short version of the project. It has three parts: the demand estimates
(Section~\ref{sec:demand}), what they imply about vertical conduct
(Section~\ref{sec:conduct}), and a counterfactual that removes random-weight private-label
cheese (Section~\ref{sec:cf}). Details, robustness checks and the specification history are in
\code{demand\_estimation\_nevo\_style.md}.

\paragraph{Main findings.}
\begin{enumerate}
  \item BLP instruments are weak in this panel (first-stage partial $R^2$ of 0.006) and their
        estimates move with the sample and the instrument subset (price coefficient $-16$ to
        $-30$). Adding wholesale cost gives a stable price coefficient of $-8.2$ (0.45) and a
        median own-price elasticity of $-1.5$. Promotions shift demand strongly (bonus buy
        0.46, price reduction 0.98 in utility units).
  \item Among the conduct models, product-by-product retail pricing fits the observed retail
        margins best. Category-wide retail pricing, manufacturer price setting and double
        marginalization are rejected; double marginalization implies negative manufacturing
        cost for 77\% of observations. Every model over-predicts the observed 34\% margin.
  \item Removing store-cut private-label cheese (14\% of volume) lowers cheese volume by
        about 10\%, consumer surplus by about 11\% and the retailer's gross margin by
        11--13\%. Only 27--31\% of the lost volume moves to packaged cheese, and remaining
        prices rise by less than 1\%.
\end{enumerate}

% ============================================================================
\section{Setup}\label{sec:setup}

\begin{table}[h!]
\centering\small
\begin{tabularx}{\textwidth}{@{}lL@{}}
\toprule
Item & Choice \\
\midrule
Data & Dominick's cheese category, 83 stores with demographics, 380 weeks \\
Market & store $\times$ week: 27,230 markets, 1,087,156 product-market observations \\
Products & 41 brand group $\times$ type group $\times$ size class cells (90\% of units) \\
Shares & 8-oz equivalent units / market size; market size $= 1.5\,\times$ the store's maximum weekly sales \\
Price & \$ per ounce, volume-weighted over the product's UPCs \\
Linear part & price, 4 size, 7 brand and 8 type dummies; store and week fixed effects absorbed \\
Random coefficients & constant, price, family-size package (diagonal $\Sigma$); constant $\times$ income, price $\times$ income, family size $\times$ household size \\
Consumers & 20 per store: log income from the store's lognormal distribution, household size at the store mean \\
BLP instruments & sums of size-class and type-group dummies over the firm's other products and over rivals' products (28); store demographics $\times$ characteristics (20) \\
Cost instruments & own wholesale cost; mean wholesale cost of rival products in the market and in the type group (3) \\
Promotion controls & fixed-weight share of the product's UPCs on bonus buy (B) and on simple price reduction (S) \\
\bottomrule
\end{tabularx}
\caption{Specification.}\label{tab:setup}
\end{table}

Code: \code{code/03\_blp\_nevo.py} (demand), \code{code/04\_supply.py}
(Sections~\ref{sec:conduct} and~\ref{sec:cf}).

% ============================================================================
\section{Demand estimates}\label{sec:demand}

Table~\ref{tab:demand} reports three full-sample estimates. Columns (1) and (2) differ only
in the instruments: (2) adds the three wholesale-cost instruments to the BLP instruments of
(1). Column (3) adds the promotion controls to (2).

\begin{table}[h!]
\centering\small
\begin{tabular}{@{}lccc@{}}
\toprule
 & (1) & (2) & (3) \\
 & BLP instruments & + wholesale cost & + wholesale cost, \\
 &  & (main) & promotion controls \\
\midrule
Price                              & $-20.90$ (0.51) & $-8.21$ (0.45) & $-6.19$ (0.46) \\
Promo B                            &                 &                & 0.46 (0.01) \\
Promo S                            &                 &                & 0.98 (0.02) \\
\addlinespace
$\sigma$ constant                  & 0.00 (0.58)     & 0.00 (0.21)    & 0.00 (0.23) \\
$\sigma$ price                     & 1.62 (0.43)     & 2.93 (0.24)    & 2.50 (0.24) \\
$\sigma$ family size               & 0.67 (0.29)     & 0.33 (0.10)    & 0.35 (0.11) \\
constant $\times$ income           & $-0.01$ (0.18)  & 0.31 (0.25)    & 0.12 (0.25) \\
price $\times$ income              & $-0.31$ (0.12)  & $-0.12$ (0.08) & $-0.08$ (0.07) \\
family size $\times$ household size & 0.34 (0.02)    & 0.36 (0.01)    & 0.36 (0.01) \\
\addlinespace
Private label                      & 0.47            & 1.19           & 1.26 \\
Kraft                              & 1.03            & 0.99           & 1.00 \\
Philadelphia                       & 1.48            & 1.59           & 1.60 \\
Sargento                           & 2.04            & 1.68           & 1.66 \\
Random weight (size class)         & 1.12            & 0.26           & 0.37 \\
\addlinespace
Own-price elasticity, median       & $-4.54$         & $-1.50$        & $-1.13$ \\
5th / 95th percentile              & $-8.33$ / $-2.92$ & $-2.41$ / $-1.00$ & $-1.83$ / $-0.75$ \\
Excluded instruments               & 46              & 49             & 49 \\
\bottomrule
\end{tabular}
\caption{Demand estimates, full sample (27,230 markets). Standard errors in parentheses.
Brand and size coefficients are relative to Borden and the bulk size class.}\label{tab:demand}
\end{table}

\subsection{BLP instruments are weak here}

\begin{table}[h!]
\centering\small
\begin{tabular}{@{}lrr@{}}
\toprule
Excluded instruments & Partial $R^2$ & $F$ \\
\midrule
BLP sums                               & 0.006 & 23 \\
Demographics $\times$ characteristics  & 0.001 & 4.6 \\
Own wholesale cost                     & 0.835 & 438,441 \\
All                                    & 0.839 & 9,775 \\
\bottomrule
\end{tabular}
\caption{First stage for price, after store and week fixed effects and the 19 dummies
(quarterly subsample, 2,170 markets).}\label{tab:firststage}
\end{table}

\begin{itemize}
  \item \textbf{Why.} Product characteristics do not vary over time, so BLP sums move only
        when the choice set changes. The 41 products are sold in almost every store-week
        (standard deviation of the number of products per market: 0.7), so the sums are close
        to constant within a product once the fixed effects and dummies are removed.
  \item \textbf{The estimates are not stable.} With the same specification the price
        coefficient is $-15.6$ (1.6) on the quarterly subsample and $-20.9$ (0.5) on the full
        sample. Dropping the demographic instruments and keeping only the BLP sums gives
        $-29.6$ (2.4), with $\sigma$ price at zero and demographic interactions of 20 and
        $-15$. With wholesale cost the two samples give $-7.5$ (1.5) and $-8.2$ (0.45).
  \item \textbf{The standard errors in column (1) should not be trusted.} They are small,
        but conventional GMM standard errors are not valid with weak instruments, and the
        instability above is the better guide to precision.
  \item \textbf{The two instrument sets disagree on the level of the elasticity}, and this
        matters for Section~\ref{sec:conduct}. Column (1) gives a median elasticity of
        $-4.5$, column (2) gives $-1.5$. The retailer's observed margin of 34\% corresponds
        to an elasticity near $-3$ under product-by-product pricing, between the two.
        Wholesale cost is far stronger, but it is not beyond doubt: it is an average
        acquisition cost, and trade deals can move it together with promotions
        (Section~\ref{sec:limits}). Column (2) is the main specification because its
        estimates are stable and its first stage is strong, not because column (1) is shown
        to be wrong.
\end{itemize}

\subsection{Promotion controls}

\begin{itemize}
  \item Both promotion coefficients are large and precisely estimated. A simple price
        reduction tag is worth about twice a bonus buy.
  \item Adding the controls lowers the price coefficient in absolute value ($-8.2$ to
        $-6.2$), and the median elasticity from $-1.5$ to $-1.1$: part of what looked like a
        price response is a response to the promotion signal.
  \item Column (2) stays the main specification for the supply side. In column (3) about
        30\% of product-markets have inelastic demand, and promotions are a retailer choice
        that the supply models below do not describe.
\end{itemize}

\subsection{Other remarks}

\begin{itemize}
  \item \textbf{$\sigma$ constant} is at its lower bound of zero in all columns: store and
        week fixed effects leave little variation in the inside share to identify it.
  \item \textbf{Optimal instruments} \citep{chamberlain1987,reynaert2014} computed from
        column (2) halve the standard errors but give a price coefficient of $-4.6$ and a
        median elasticity of $-0.82$, which no profit-maximizing conduct can rationalize.
        They are not used; see \code{demand\_estimation\_nevo\_style.md}.
\end{itemize}

% ============================================================================
\section{Vertical conduct}\label{sec:conduct}

\subsection{Approach}

Each store-week has one retailer, and Dominick's records its wholesale cost $w$ for every
UPC. So the retail margin $p - w$ is observed (median 34\% of price), and the retail side of
each conduct model can be compared with data directly. The models follow
\citet{villasboas2007}.

\begin{table}[h!]
\centering\small
\begin{tabularx}{\textwidth}{@{}lLL@{}}
\toprule
Model & Who sets what & Total margin $p - c$ \\
\midrule
\model{R-prod}  & Retailer prices each product on its own, given $w$ &
                  $-s_j / (\partial s_j / \partial p_j)$ \\
\model{R-brand} & Retailer prices brand by brand, given $w$ &
                  Bertrand, ownership $=$ brand group \\
\model{R-cat}   & Retailer maximizes category profit, given $w$. With nonlinear wholesale
                  contracts ($w =$ manufacturer marginal cost) this is the vertically
                  integrated outcome &
                  $-(\partial s / \partial p')^{-1} s$ \\
\model{M-bert}  & Manufacturers set retail prices; retail margin is zero; the retailer owns
                  the private label &
                  Bertrand, ownership $=$ manufacturer \\
\model{DM}      & Double marginalization: manufacturers set linear wholesale prices, retailer
                  then maximizes category profit &
                  \model{R-cat} retail margin $+$ manufacturer margin through retail
                  pass-through \\
\bottomrule
\end{tabularx}
\caption{Conduct models.}\label{tab:models}
\end{table}

For \model{R-prod}, \model{R-brand} and \model{R-cat} the model margin is a prediction of
the observed $p - w$. Under \model{M-bert} the predicted retail margin is zero. Under
\model{DM} the retail part is the \model{R-cat} margin, and the manufacturer margin is
\[
  w - c = -\bigl(T_w \circ (\Delta_p G)'\bigr)^{-1} s ,
\]
where $\Delta_p = \partial s / \partial p'$, $T_w$ is the manufacturer ownership matrix and
$G = \partial p / \partial w'$ is the pass-through matrix from the retailer's first-order
conditions (it uses the second derivatives of demand). Private-label wholesale prices are
treated as non-strategic.

Mean utilities are recovered by share inversion at the main estimates (column 2 of
Table~\ref{tab:demand}); the implied elasticities match pyBLP's to eight digits.

\subsection{Results}

\begin{table}[h!]
\centering\small
\begin{tabular}{@{}lcccc@{}}
\toprule
 & Median & Quantity-weighted & Implied & RMSE vs observed \\
 & margin / $p$ & margin / $p$ & $mc < 0$ & retail margin (share of $p$) \\
\midrule
Observed retail margin & 0.34 & 0.30 &       &      \\
\model{R-prod}         & 0.67 & 0.83 & 5\%   & 0.41 \\
\model{R-brand}        & 0.71 & 0.88 & 8\%   & 0.46 \\
\model{M-bert}         & 0.73 & 0.93 & 12\%  & 0.49 \\
\model{R-cat}          & 0.92 & 1.17 & 38\%  & 0.69 \\
\model{DM}             & 1.26 & 1.70 & 77\%  &      \\
\bottomrule
\end{tabular}
\caption{Implied total margins by conduct model, full sample.}\label{tab:margins}
\end{table}

\begin{table}[h!]
\centering\small
\begin{tabular}{@{}cccc@{}}
\toprule
\model{R-prod} vs \model{R-brand} & \model{R-prod} vs \model{R-cat} &
\model{R-brand} vs \model{R-cat} & \model{M-bert} vs \model{R-prod} \\
\midrule
$t = -25.5$ & $t = -25.2$ & $t = -24.9$ & $t = +28.3$ \\
\bottomrule
\end{tabular}
\caption{Pairwise comparison of fit to the observed retail margin: difference in squared
errors, standard errors clustered by store. Negative $t$: the first model fits
better.}\label{tab:pairwise}
\end{table}

Under \model{R-cat} the median pass-through of own wholesale cost to the retail price is
1.25, close to the 1.27 in the raw data. The \model{DM} manufacturer margin on national
brands is 54\% of the retail price at the median.

\begin{table}[h!]
\centering\small
\begin{tabular}{@{}lccc@{}}
\toprule
 & National & Private label, & Private label, \\
 & brands & packaged & random weight \\
\midrule
Share of volume               & 61\%    & 25\%    & 14\% \\
Price, \$/oz                  & 0.204   & 0.174   & 0.187 \\
Observed retail margin / $p$  & 0.24    & 0.37    & 0.40 \\
Median own elasticity         & $-1.59$ & $-1.39$ & $-1.30$ \\
\model{R-prod} margin / $p$   & 0.82    & 0.88    & 0.78 \\
\model{R-cat} margin / $p$    & 1.16    & 1.24    & 1.11 \\
\bottomrule
\end{tabular}
\caption{Margins by product group (quantity-weighted).}\label{tab:groups}
\end{table}

\subsection{What the test says}

\begin{itemize}
  \item \textbf{The model with the least internalization fits best.} Implied margins rise in
        the order \model{R-prod} $<$ \model{R-brand} $<$ \model{M-bert} $<$ \model{R-cat}
        $<$ \model{DM}, and all exceed the observed 34\%. A \citet{millerweinberg2017} style
        internalization parameter $\kappa$ (ownership matrix with $\kappa$ off the diagonal)
        would be estimated at its lower bound of zero.
  \item \textbf{Double marginalization is rejected.} It needs a manufacturer margin of 54\%
        of the retail price on top of the retail margin, which implies negative
        manufacturing cost for 77\% of observations.
  \item \textbf{Category-wide retail pricing is rejected} against product-by-product
        pricing: it implies margins near 92\% of price and negative cost in 38\% of
        observations. Dominick's behaves like a retailer that takes wholesale prices as
        given and does not fully exploit category-level market power.
  \item \textbf{Manufacturer retail-price setting is rejected directly}: it predicts a zero
        retail margin, and the observed margin is 34\% (24\% on national brands, 37--40\% on
        private label).
  \item \textbf{Even the best model over-predicts}: \model{R-prod} implies 67\% against 34\%
        observed, which would need a median own elasticity near $-3$. Likely reasons:
        products are coarse aggregates (substitution within a cell is not counted), weekly
        data mix promotion timing with long-run response, and competition from other chains
        is outside a store-level market. The private-label premium in observed margins
        (37--40\% vs 24\%) is not reproduced either, which points to wholesale prices of
        national brands being above marginal cost.
\end{itemize}

Conduct used in Section~\ref{sec:cf}: \model{R-prod}.

% ============================================================================
\section{Counterfactual: packaged cheese only}\label{sec:cf}

\subsection{Design}

Suppose store-cut, random-weight cheese can no longer be sold (for example for hygiene
reasons). In the data this is four private-label products: cheddar, other natural,
mozzarella and processed, sold by weight. They are present in every market and account for
14.0\% of inside volume, at an observed margin of 40\%.

The four products are removed from every choice set. Two versions:
\begin{itemize}
  \item \textbf{Prices fixed.} Uses demand only.
  \item \textbf{Re-priced.} Remaining prices are re-solved under \model{R-prod}, holding
        fixed the marginal cost that \model{R-prod} implies at observed prices.
\end{itemize}

Consumer surplus is the expected maximum utility divided by each consumer's marginal utility
of a dollar. Retail gross margin is $(p - w)\,q$ at observed wholesale costs. Baseline per
store-week: revenue \$6,216, gross margin \$1,948.

\subsection{Results}

\begin{table}[h!]
\centering\small
\begin{tabular}{@{}lcc@{}}
\toprule
 & Prices fixed & Re-priced (\model{R-prod}) \\
\midrule
Cheese volume                               & $-9.7\%$  & $-10.3\%$ \\
Lost volume diverted to remaining cheese    & 31\%      & 27\% \\
\addlinespace
Remaining private label: price              & 0         & $+0.7\%$ \\
Remaining national brands: price            & 0         & $+0.9\%$ \\
Remaining private label: quantity           & $+4.9\%$  & $+4.3\%$ \\
Remaining national brands: quantity         & $+5.1\%$  & $+4.3\%$ \\
\addlinespace
Consumer surplus                            & $-10.7\%$ & $-11.2\%$ \\
Consumer surplus, \$ per store-week         & $-869$    & $-916$ \\
Retail gross margin                         & $-12.8\%$ & $-11.2\%$ \\
Retail gross margin, \$ per store-week      & $-249$    & $-218$ \\
Variable profit at \model{R-prod} costs     & $-9.3\%$  & $-9.1\%$ \\
\bottomrule
\end{tabular}
\caption{Removing random-weight private-label cheese. The pricing equilibrium was found in
95.2\% of markets; the re-priced column uses those.}\label{tab:cf}
\end{table}

\subsection{Interpretation}

\begin{itemize}
  \item \textbf{Volume.} About 70\% of the removed volume leaves the category. Substitution
        to packaged cheese is spread almost evenly: packaged private label and national
        brands both gain 4--5\%, and packaged products of the same types gain no more than
        other types. This is the logit pattern. The random coefficients are on the constant
        (at zero), price and family size, so the model has no dimension along which
        store-cut cheddar is closer to packaged cheddar than to cream cheese. Diversion to
        close substitutes is probably understated and the volume loss overstated.
  \item \textbf{Prices.} Remaining prices rise by less than 1\%. The effect is small because
        under \model{R-prod} the retailer did not internalize substitution between products
        to begin with; removing products only makes the remaining demand slightly less
        elastic.
  \item \textbf{Consumers.} Surplus falls by about 11\%, or roughly \$900 per store-week
        against category revenue of \$6,216. The figure includes the loss of the logit taste
        draw for four products, which tends to overstate the value of variety.
  \item \textbf{Retailer.} Gross margin falls by 13\% at fixed prices: the removed products
        carry the highest margins (40\%), and the diverted volume goes partly to national
        brands at 24\%. Re-pricing recovers about one eighth of the loss.
  \item \textbf{Manufacturers.} National-brand volume rises by 4--5\%, so branded
        manufacturers gain. Their profit change is not reported in dollars because their
        costs are not identified without a conduct model that fits.
\end{itemize}

% ============================================================================
\section{Limitations}\label{sec:limits}

\begin{itemize}
  \item All conduct models over-predict the retail margin, so the marginal costs used in the
        re-priced counterfactual are too low. The fixed-price column does not depend on
        this.
  \item Wholesale cost in Dominick's data is the average acquisition cost of inventory, not
        the replacement cost, so observed margins lag wholesale price changes.
  \item The conduct comparison is on retail margins only. Manufacturer marginal costs are
        not observed, so the upstream part of each model is judged by whether implied costs
        are positive, not by a formal cost-side test as in \citet{villasboas2007}.
  \item The counterfactual holds wholesale prices, the product set of packaged cheese and
        store traffic fixed.
\end{itemize}

% ============================================================================
\begin{thebibliography}{9}

\bibitem[Chamberlain(1987)]{chamberlain1987}
Chamberlain, G. (1987). Asymptotic efficiency in estimation with conditional moment
restrictions. \emph{Journal of Econometrics}, 34(3), 305--334.

\bibitem[Conlon and Gortmaker(2020)]{conlon2020}
Conlon, C. and Gortmaker, J. (2020). Best practices for differentiated products demand
estimation with PyBLP. \emph{RAND Journal of Economics}, 51(4), 1108--1161.

\bibitem[Miller and Weinberg(2017)]{millerweinberg2017}
Miller, N. H. and Weinberg, M. C. (2017). Understanding the price effects of the MillerCoors
joint venture. \emph{Econometrica}, 85(6), 1763--1791.

\bibitem[Nevo(2000)]{nevo2000}
Nevo, A. (2000). A practitioner's guide to estimation of random-coefficients logit models of
demand. \emph{Journal of Economics \& Management Strategy}, 9(4), 513--548.

\bibitem[Reynaert and Verboven(2014)]{reynaert2014}
Reynaert, M. and Verboven, F. (2014). Improving the performance of random coefficients demand
models: the role of optimal instruments. \emph{Journal of Econometrics}, 179(1), 83--98.

\bibitem[Villas-Boas(2007)]{villasboas2007}
Villas-Boas, S. B. (2007). Vertical relationships between manufacturers and retailers:
inference with limited data. \emph{Review of Economic Studies}, 74(2), 625--652.

\end{thebibliography}

\end{document}
